%% Chapter 2 — Related Work
\chapter{Related Work}
\label{ch:related_work}

\section{Non-Photorealistic Rendering for 3D Characters}
\label{sec:related_npr}

Non-photorealistic rendering has a substantial history in computer graphics, though most of it predates the era of closed, SDK-managed avatar systems. Understanding this body of work requires distinguishing between the underlying techniques — edge detection, shading stylisation, and painterly filtering — and their specific application to animated human figures.

The foundational contribution to NPR lighting is the cool-to-warm shading model introduced by Gooch et al.~\cite{gooch1998}, which replaced the continuous Lambertian diffuse term with a hue shift from cool shadows to warm highlights, mimicking the conventions of technical illustration. While the model targets engineering drawings rather than characters, its core insight — that shading can communicate form without physically simulating light — became a reference point for subsequent stylised rendering work. In the context of this thesis, the more direct technical heritage lies in outline extraction and edge-based effects rather than shading models.

Silhouette rendering has been approached from two geometrically distinct directions. Geometry-based methods generate outlines using the GPU rasteriser itself. The inverted hull technique, described by Lake et al.~\cite{lake2000}, expands back-facing polygons along vertex normals by a small offset, creating a slightly enlarged shell of the mesh. When the front-facing geometry is rendered on top, only the protruding rim remains visible as an outline. The approach requires no image-space computation, scales predictably with polygon count, and remains consistent regardless of camera distance — properties that make it well suited to animated characters in real-time applications. Isenberg et al.~\cite{isenberg2003} provide a systematic survey of silhouette algorithms and conclude that geometry-based methods are robust and distance-independent, but structurally incapable of detecting internal surface detail. This limitation matters for avatar rendering, where facial features, clothing folds, and skin creases all constitute expressive visual cues that lie within the silhouette boundary.

Image-space methods operate on the rasterised output rather than on the mesh directly. Screen-space normal derivatives, computed via the hardware-accelerated \texttt{ddx}/\texttt{ddy} intrinsics available on all modern GPUs, detect abrupt changes in surface orientation without any additional draw calls or intermediate render targets. This approach captures surface creases and mesh boundaries that geometry expansion misses. The fundamental limitation is that edge thickness is measured in screen pixels rather than world-space units, so edges scale with camera distance. In an immersive VR context where the user moves relative to other participants, this produces edges that appear overly thick at close range and invisible at distance.

Texture-based edge detection targets a different layer of information. The Sobel operator~\cite{gonzalez2018}, a discrete gradient approximation applied to the diffuse texture via a $3\times3$ convolution kernel, identifies luminance boundaries that correspond to clothing patterns, facial contours, and painted surface details. Unlike geometry-based methods, it responds to what was painted onto the surface rather than to the surface shape itself. The approach is well established in image processing~\cite{canny1986}, but differentiation inherently amplifies high-frequency noise. Marr and Hildreth~\cite{marr1980} demonstrated that low-pass pre-filtering is a prerequisite for stable gradient computation, an observation that motivates the Gaussian pre-filtered variant explored later in this work.

Extended toon shading approaches have explored richer parameter spaces. Barla et al.~\cite{barla2006} introduced X-Toon, which parameterises the shading tone ramp by both view angle and depth, enabling a two-dimensional lookup that encodes artistic intent unavailable in a one-dimensional quantisation step. On the painterly side, the Kuwahara filter — originally a noise-reduction tool in medical imaging — has been applied to NPR as a means of producing region-based painterly abstraction without explicit edge detection. Kyprianidis et al.~\cite{kyprianidis2009} extended the isotropic variant with an anisotropic formulation aligned to local image structure, producing brush-stroke effects oriented along salient contours. The anisotropic extension requires a structure tensor computed from a prior render pass and cannot be implemented within a single-pass fragment hook — a constraint directly relevant to the integration target in this thesis.

Hatching represents a separate branch of character NPR in which tonal regions are conveyed by oriented stroke textures rather than flat fills or edge lines~\cite{praun2001}. While visually distinctive, the approach depends heavily on stroke coherence across animation frames, a problem that has not been fully solved in real-time settings. It was not pursued in this work.

What connects nearly all of these contributions is an implicit assumption: the developer controls the rendering pipeline. Techniques are designed and evaluated against custom meshes, open rendering frameworks, or research engines where shader code, render passes, and material systems can be freely modified. This assumption does not hold for production avatar SDKs.

\section{Avatar Rendering and Appearance in VR}
\label{sec:related_vr_avatars}

Commercial avatar platforms have matured considerably beyond the research systems used in early social VR studies. The Meta Avatars SDK, used in Horizon Worlds and accessible to third-party developers, provides production-quality physically-based rendering with subsurface scattering for skin, procedural eye glints, hair rendering, and compute-based GPU skinning driven by body tracking data. The pipeline is generated and managed by Meta's Extensible Shader Framework~(ESF), which exposes a small number of application-specific fragment hooks while keeping the core lighting computations proprietary and unmodifiable. This architecture maintains pipeline integrity across SDK updates and ensures consistent quality across hardware, but it substantially constrains where custom rendering logic can be inserted.

Prior NPR work applied to avatar systems has generally assumed open pipeline access. Studies applying toon shading or stylisation to VR avatars have used custom character models with fully accessible shader code, or have modified appearance at the asset level before the render pipeline is involved. Applying NPR to a closed, compute-skinned, LOD-managed avatar at runtime — without SDK modification, without additional render passes, and without breaking the material management system — has not been addressed in the published literature to the author's knowledge.

The most directly related work on avatar stylisation and perception is Canales et al.~\cite{canales2024}, who conducted a controlled study of how avatar stylisation level affects perceived trust. Participants interacted with avatars at varying degrees of abstraction, from photorealistic to explicitly cartoon-like, and rated trustworthiness using social cognition instruments. The study found that moderate stylisation did not reduce trust relative to photorealistic rendering, and could in some conditions improve it, consistent with the uncanny valley framework. However, the stylisation was applied to custom assets in an open pipeline, and the technique variations compared were coarse-grained — the distinction between different edge detection approaches, painterly filtering methods, or composite techniques was not investigated. The work demonstrates that abstraction matters; it leaves open the question of which specific technical approach is most effective and whether these findings generalise to SDK-managed avatars.

Weidner and Broll~\cite{weidner2023} examined how avatar visual fidelity interacts with self-embodiment and social perception in video conferencing and VR contexts. Their findings indicate that the visual design of an avatar does not operate independently of the observer's sense of identification with it: fidelity effects on social judgements are mediated by embodiment. This has methodological implications for user studies that compare avatar conditions — a participant's prior exposure to and familiarity with a particular avatar style can influence their trust ratings independently of the stylisation itself.

\section{Trust, Social Presence, and the Uncanny Valley}
\label{sec:related_vr_trust}

The uncanny valley hypothesis, first articulated by Mori~\cite{mori1970}, proposes that human affinity for a human-like representation increases with resemblance up to a point, then drops sharply before recovering near perfect replication. The drop is attributed to a mismatch: a near-realistic representation triggers expectations calibrated for real humans, but subtle violations — in motion timing, skin texture, eye movement — produce unease rather than familiarity. MacDorman and Ishiguro~\cite{macdorman2006} provided empirical support for this in the context of humanoid robots, documenting the eeriness ratings that emerge when a representation is close but not convincingly human.

Modern avatar systems sit squarely in the region of this curve that is most problematic. They achieve surface realism sufficient to trigger high perceptual expectations, but cannot fully meet those expectations in motion timing, facial micro-expressions, or gaze behaviour. The practical consequence for social VR is that a photorealistic avatar may actively undermine the sense of trustworthiness that it was designed to support. NPR stylisation provides a principled alternative: by moving the representation away from near-photorealism and into a clearly artistic register — cel shading, ink-line treatment, painterly filtering — the avatar exits the uncanny valley zone rather than attempting to cross it.

The relationship between avatar appearance and social presence has been studied extensively. Nowak and Biocca~\cite{nowak2003} demonstrated that anthropomorphism and perceived agency are distinct constructs that contribute independently to social presence and copresence. Counterintuitively, very high levels of visual anthropomorphism did not consistently produce higher social presence scores. This suggests that the social and perceptual response to an avatar is not a simple monotonic function of visual realism, and that stylisation does not necessarily carry a social cost.

Trust in the context of virtual agents is related to social presence but distinct from it. A user can experience strong co-presence with an avatar while trusting the entity it represents very little, or trust a stylised character precisely because it does not trigger uncanny valley discomfort. The literature on avatar trust has largely compared coarse conditions — human versus clearly robotic, photorealistic versus cartoonish — without exploring the design space within NPR itself. Whether a Kuwahara-painterly avatar is trusted differently from a Sobel-edge-line avatar, or whether hierarchical multi-cue edge detection produces a better trust outcome than a simpler screen-space method, are questions that have not been addressed.

\section{Summary and Positioning}
\label{sec:related_positioning}

The technical NPR literature provides a rich set of established techniques, developed and validated in open rendering environments. The social VR literature demonstrates that avatar appearance significantly affects trust and social presence, and that the uncanny valley poses a genuine design challenge for photorealistic avatar systems. Canales et al. showed that stylisation can improve trust outcomes; Nowak and Biocca established that the appearance–presence relationship is non-linear. What is missing is the intersection: a systematic investigation of fine-grained NPR technique choice on a production, SDK-managed avatar system, paired with an empirical measure of the trust response those choices produce.

This thesis addresses that gap. It integrates a suite of NPR techniques into the Meta Avatars SDK under real production constraints, evaluates the visual and performance trade-offs across techniques, and measures the trust implications of different abstraction levels in a user study.
